\chapter{Appendix: The \texorpdfstring{$\hbar$}{hbar} parameters and the q-convention bridge}
\label{app:q-conventions-v2}

The Vol~I q-convention bridge fixes the comparison between the
Kazhdan--Lusztig and Drinfeld--Kohno multiplicative parameters.

\begin{remark}[Cross-volume q-convention reference]
\label{rem:q-bridge-vol1-canonical-pointer}
The q-convention bridge is the fully expanded Vol~I appendix at
Appendix~\ref{app:q-conventions}. Its canonical theorem has six
clauses: squaring, logarithm, the $r$-matrix bridge with trace-form
level-prefix verification at $k=0$, $R$-matrix monodromy, conformal
weights at q-deformation, and categorical agreement. The
Chriss--Ginzburg form states the bridge as the index-$2$
$\mathbb{Z}/2$-cover $B_2 \to S_2$ of the Artin braid group, realised
algebraically as the squaring $q_{\KL}^2 = q_{\DK}$. The present Vol~II
appendix supplies the nine $\hbar$-parameters (three equivalence classes A, B, C) and
the BCFG normalisation; the universal six-clause bridge is Vol~I.
\end{remark}
 That theorem
records the $\mathbb{Z}/2$-cover relation $q_{KL}^{2}=q_{DK}$ between the
Kazhdan--Lusztig and Drinfeld--Kohno multiplicative parameters attached to a
simple Lie algebra $\mathfrak{g}$ at non-critical level $k$. The present
appendix catalogues the nine formally distinct $\hbar$'s, each with
its proper domain and its map to the reference parameter
$\hbar_{\mathrm{ref}} := 1/(k+h^\vee)$, and fixes the BCFG
normalisation where the level parameter and the Drinfeld logarithm
would otherwise be confused.

\section{The nine $\hbar$'s of Vol II}
\label{sec:hbar-zoo-vol2}

Each $\hbar$ is declared with: (a) name and symbol, (b) defining formula, (c)
domain of use, (d) bridge to $\hbar_{\mathrm{ref}} := 1/(k+h^\vee)$.

\begin{enumerate}
\item[(1)] \emph{KZ parameter} $\hbar_{\mathrm{KZ}} := 1/(k+h^\vee) = \hbar_{\mathrm{ref}}$.
 The coefficient of the KZ r-matrix $r_{KZ}(z) = \hbar_{\mathrm{KZ}}\,\Omega/z$.
 Domain: affine half-space BV, spectral braiding core, log HT monodromy.
 Bridge: identity.
\item[(2)] \emph{Sugawara parameter} $\hbar_{\mathrm{Sug}} := 1/(2(k+h^\vee))$.
 The coefficient in $T(z) = \hbar_{\mathrm{Sug}} \sum_a {:}J^a J_a{:}(z)$.
 Domain: topologisation and holographic reconstruction.
 Bridge: $\hbar_{\mathrm{Sug}} = \hbar_{\mathrm{ref}}/2$.
\item[(3)] \emph{Yangian deformation parameter} $\hbar_{\mathrm{Yan}}$.
 The formal parameter of $Y_{\hbar_{\mathrm{Yan}}}(\mathfrak{g})$, independent
 of any level. Declared as an indeterminate over the ground ring.
 Domain: gravitational Yangian and type-A Baxter--Rees material.
 Bridge: $\hbar_{\mathrm{Yan}}$ is identified with $\hbar_{\mathrm{ref}}$ only
 after rational degeneration of the affine KM at level $k$; outside that
 limit, independent.
\item[(4)] \emph{Drinfeld quantisation parameter} $\hbar_{\mathrm{Dr}}$.
 Defined by $q_{\mathrm{D}} = \exp(\hbar_{\mathrm{Dr}})$, so
 $\hbar_{\mathrm{Dr}} = \log q_{\mathrm{D}}$.
 Domain: spectral Drinfeld strictification, DNP identification.
 Bridge: $\hbar_{\mathrm{Dr}} = 2\pi i\,\hbar_{\mathrm{ref}}$ (principal branch),
 so $q_{\mathrm{D}} = q_{DK}$.
\item[(5)] \emph{Kazhdan parameter} $\hbar_{\mathrm{K}}$.
 Defined by $q_{\mathrm{K}} = \exp(\pi i\,\hbar_{\mathrm{K}})$, i.e.\ the
 half-monodromy parameter.
 Domain: wherever $q_{KL}$ is used as $q_{\mathrm{K}}$, including the Vol I
 bridge reference.
 Bridge: $\hbar_{\mathrm{K}} = \hbar_{\mathrm{ref}}$, $q_{\mathrm{K}} = q_{KL}$.
\item[(6)] \emph{BV homological parameter} $\hbar_{\mathrm{BV}}$.
 The formal variable deforming the classical BV master equation to the
 quantum master equation $(\mathrm{d}+\{-,S\}+\hbar_{\mathrm{BV}}\Delta_{\mathrm{BV}})^{2}=0$.
 Domain: BV--BRST, affine half-space BV, and HT physics.
 Bridge: $\hbar_{\mathrm{BV}}$ is a \emph{distinct} indeterminate; its
 identification with $\hbar_{\mathrm{ref}}$ holds only after the factorisation
 algebra of CFG is restricted to affine KM, where
 $\hbar_{\mathrm{BV}} = \hbar_{\mathrm{ref}}$ is the content of the
 Costello--Li reduction.
\item[(7)] \emph{Brace/operadic parameter} $\hbar_{\mathrm{br}}$.
 Genus expansion parameter in the brace / genus tower
 $\sum_{g\ge 0} \hbar_{\mathrm{br}}^{\,g}\,F_g$.
 Domain: brace operations, chiral Hochschild theory, and modular bootstrap.
 Bridge: $\hbar_{\mathrm{br}}$ is genus-weight, unrelated to level. No bridge.
\item[(8)] \emph{Celestial / soft parameter} $\hbar_{\mathrm{cel}}$.
 The deformation parameter in the soft graviton tower
 tracking the celestial
 $w_{1+\infty}$ expansion.
 Bridge: $\hbar_{\mathrm{cel}}$ equals $\hbar_{\mathrm{br}}$ at genus $0$ and
 differs beyond; not equal to $\hbar_{\mathrm{ref}}$.
\item[(9)] \emph{Omega-background parameter} $\hbar_{\Omega}$ (and its
 refinement $(\epsilon_{1},\epsilon_{2})$ with $\hbar_{\Omega}=\epsilon_{1}+\epsilon_{2}$).
 Domain: BCFG folding, gauge-theoretic Yangian origins,
 Yang--Mills synthesis and instanton screening.
 Bridge: $\hbar_{\Omega}$ is an equivariant parameter on the instanton moduli
 space, independent of $\hbar_{\mathrm{ref}}$; the Nekrasov/AGT bridge
 identifies $\hbar_{\Omega}$ with $\hbar_{\mathrm{Yan}}$ after the gauge-theory
 limit.
\end{enumerate}

\begin{remark}[Three equivalence classes]
The nine $\hbar$'s fall into three equivalence classes under the Vol I bridge:
(A) \emph{Level class}: $\hbar_{\mathrm{KZ}},\hbar_{\mathrm{Sug}},\hbar_{\mathrm{K}}$
(all proportional to $\hbar_{\mathrm{ref}}$). (B) \emph{Quantisation class}:
$\hbar_{\mathrm{Dr}},\hbar_{\mathrm{Yan}},\hbar_{\mathrm{BV}},\hbar_{\Omega}$
(formal indeterminates, paired to class (A) by specific bridge theorems). (C)
\emph{Grading class}: $\hbar_{\mathrm{br}},\hbar_{\mathrm{cel}}$ (genus /
soft-expansion weights, unrelated to level). Conflating classes is the source
of the cascade B-cycle $i^{2}$ sign error and the BCFG
normalisation clash below.
\end{remark}

\section{BCFG \texorpdfstring{$\hbar$}{hbar} normalisation}
\label{sec:bcfg-hbar-resolution}

BCFG folding admits two incompatible declarations of a symbol $\hbar$: one as
$\hbar = 1/(k+2)$ (an instance of $\hbar_{\mathrm{ref}}$ for $\mathfrak{sl}_{2}$,
i.e.\ class (A)), and one as $\hbar = \pi i/(k+2)$ (secretly
$\hbar = \hbar_{\mathrm{Dr}}/2$, class (B)). Under the second declaration the
B-cycle monodromy computation produces
$q = e^{2\pi i\hbar} = e^{2\pi i\cdot\pi i/(k+2)} = e^{-2\pi^{2}/(k+2)}$, which
is \emph{real} and less than unity: the $i^{2}=-1$ silently turns a root of
unity into a real exponential decay. This is the B-cycle $i^{2}$ sign error
arising from a clash between the level parameter and the Drinfeld
logarithm.

\begin{resolution}
Throughout all descendants of the folded BCFG chapter, the symbol $\hbar$ is
fixed to mean $\hbar_{\mathrm{ref}} = 1/(k+h^\vee)$ (class (A), item (1)
above). The Drinfeld parameter is written explicitly as
$\hbar_{\mathrm{Dr}} = 2\pi i\,\hbar_{\mathrm{ref}}$ whenever needed. The
B-cycle monodromy is then
\[
 q_{DK} \;=\; \exp(\hbar_{\mathrm{Dr}}) \;=\; \exp(2\pi i\,\hbar_{\mathrm{ref}}),
\]
which lies on the unit circle and is a root of unity precisely when
$\hbar_{\mathrm{ref}}\in\mathbb{Q}$, as required by Kazhdan--Lusztig.
\end{resolution}

\section{Cross-volume convention table}
\label{sec:cross-vol-conventions}

\begin{center}
\begin{tabular}{llll}
\hline
Symbol & Vol I (canonical) & Vol II (this volume) & Relation \\
\hline
$\hbar_{\mathrm{ref}}$ & $1/(k+h^\vee)$ & $\hbar_{\mathrm{KZ}}=\hbar_{\mathrm{K}}$ & identity \\
$q_{KL}$ & $\exp(\pi i\hbar_{\mathrm{ref}})$ & $q_{\mathrm{K}}$ & identity \\
$q_{DK}$ & $\exp(2\pi i\hbar_{\mathrm{ref}})$ & $q_{\mathrm{D}}$ & $q_{DK}=q_{\mathrm{D}}$ \\
$\hbar_{\mathrm{Dr}}$ & $\log q_{DK}$ & $\log q_{\mathrm{D}}$ & $=2\pi i\,\hbar_{\mathrm{ref}}$ \\
$\hbar_{\mathrm{Yan}}$ & Yangian deformation & Yangian deformation & $=\hbar_{\mathrm{ref}}$ in rational limit \\
$\hbar_{\mathrm{BV}}$ & CFG factorisation weight & BV homological weight & $=\hbar_{\mathrm{ref}}$ after Costello--Li \\
$\hbar_{\mathrm{Sug}}$ & absent & $1/(2(k+h^\vee))$ & $=\hbar_{\mathrm{ref}}/2$ \\
$\hbar_{\mathrm{br}}$ & genus weight & genus weight & independent \\
$\hbar_{\mathrm{cel}}$ & absent & soft-expansion weight & $=\hbar_{\mathrm{br}}$ at $g{=}0$ \\
$\hbar_{\Omega}$ & absent & equivariant parameter & $=\hbar_{\mathrm{Yan}}$ in gauge limit \\
\hline
\end{tabular}
\end{center}

\begin{theorem}[Vol II consistency with the Vol I bridge]
\label{thm:q-convention-bridge-v2}
Under the identifications of the table above, the six assertions
\textup{(i)--(vi)} of the q-convention bridge hold verbatim in
Vol II with $\hbar = \hbar_{\mathrm{ref}} = \hbar_{\mathrm{KZ}} = \hbar_{\mathrm{K}}$.
In particular, for the ordered bar complex $B^{\mathrm{ord}}(\widehat{\mathfrak{g}}_{k})$
on two adjacent chiral inputs, the braiding is $q_{KL}^{\Omega}$
(half-monodromy on the $\mathbb{Z}/2$-cover) and the class of a pure braid is
$q_{DK}^{\Omega} = q_{KL}^{2\Omega}$ (full monodromy on the base).
\end{theorem}

\begin{proof}
(i)--(vi) are the Vol~I q-convention bridge with the table
substitutions. The bar-complex statement is the specialisation of assertion
(iv) to the E$_{1}$-chiral bar complex: the adjacent-swap generator of the
Artin braid group on ordered configurations of two marked points in
$\mathrm{Conf}_{2}(\mathbb{C})$ acts as $q_{KL}^{\Omega}$ on the evaluation
module $V_{\lambda}\boxtimes V_{\mu}$, and the pure-braid generator as its
square.
\end{proof}

\begin{remark}[Failure modes]
The Vol~II pitfalls around $\hbar$ and $q$ neutralised here are: the
B-cycle $i^{2}$ sign error (a root of unity silently
replaced by an exponential decay); $\hbar$ convention clashes; the
misconstruction of the universal R-matrix from $\Delta(1)$ in place of the
half-braiding presentation; the Verdier-dual level identification $k^{!} = -k$
and its interaction with the $\sigma_{2}$-parity of the Shapovalov form; and
the distinction between the Weierstrass $\wp_{1}$ propagator (periodic) and the
$\theta_{1}'/\theta_{1}$ propagator (quasi-periodic with distinct monodromy
data). Section~\ref{sec:hbar-zoo-vol2} supplies a named symbol for every
occurrence so that none of these conflations can silently recur.
\end{remark}
